\subsection{Analyse}

\paragraph{Cohérence avec l'angle $A$ du prisme} 
La valeur mesurée de l'angle du prisme ($A = (60.3 \pm 0.5)^{\circ})$
est en accord avec les caractéristiques du prisme utilisé pour les manipulations. 

\paragraph{Nature du matériau}
L'indice de réfraction calculé ($n \approx 1.733$) correspondent à celui d'un verre dense type Flint, 
cohérent avec les indications reçues en séance. 

\paragraph{Analyse de l'incertitude}
Malgré l'incertitude faible, il est à noter que les mesures auraient pu gagner en précision. 
En effet, les mesures en minutes d'arc étaient impossible à cause d'un défaut du goniomètre. 

\subsection{Discussion}

\paragraph{Apprentissage et méthodologie}
Nous avons eu beaucoup de difficulté au départ à être certains de nos mesures. 
En effet, nous n'étions absolument pas certains de ce que nous devions observer au départ, 
et nous avons eu des difficultés à évaluer si une mesure était cohérente 
d'un point de vue physique/théorique. 
Nous nous sommes rendus compte également d'à quel point nous avions une marge de progression
quand nous nous sommes mis à aider d'autres groupes de TP passés après nous. 
Aujourd'hui, si nous devions recommencer l'expérience, 
il nous parait assez clair que certains éléments de celle-ci se passerait mieux. 

